\documentclass[11pt]{article}
\usepackage[T1]{fontenc}
\usepackage{lmodern,microtype,amsmath,amssymb,amsthm,booktabs,longtable,array}
\usepackage[margin=1in]{geometry}
\usepackage{xurl}
\usepackage[hidelinks]{hyperref}
\setlength{\emergencystretch}{3em}
\newtheorem{theorem}{Theorem}
\newtheorem{lemma}[theorem]{Lemma}
\newtheorem{proposition}[theorem]{Proposition}
\providecommand{\tightlist}{\setlength{\itemsep}{0pt}\setlength{\parskip}{0pt}}
\providecommand{\passthrough}[1]{#1}
\hypersetup{pdfauthor={Yongxian Zhang}}

\title{Infinite log-concavity of cycle chromatic coefficients: a complete classification}
\author{Yongxian Zhang\\[2pt]\small School of Computer Science and Engineering\\\small South China University of Technology\\\small\texttt{mxymmxym1@gmail.com}\\\small ORCID: 0009-0000-3864-3536}
\date{8 October 2026}
\begin{document}\maketitle
\begin{abstract}We classify infinite log-concavity of the full absolute chromatic coefficient sequence of every cycle: it holds exactly for lengths 3 through 11. Explicit negative iterates handle lengths 12 through 16, and a uniform third-iterate identity excludes all larger lengths. A strengthened log-concavity invariant proves all infinitely many positive iterates. The complete chain, from actual graph colorings to the if-and-only-if classification, has a Lean proof.\end{abstract}
\section{Statement and context}
Let $C_n$ be the undirected cycle with $n\ge3$ vertices. For its chromatic
polynomial $\chi_n(q)$, write $a_k=|[q^k]\chi_n(q)|$, and extend this
sequence by zero to every $k\in\mathbb Z$ outside $0\le k\le n$.
For any zero-extended finite sequence define
\[
 (\mathcal La)_k=a_k^2-a_{k-1}a_{k+1}.
\]
A sequence is infinitely log-concave if every entry of $\mathcal L^r a$
is nonnegative for every integer $r\ge0$. The endpoint positions are
retained at every iteration.

\begin{theorem}\label{thm:classification}
The absolute coefficient sequence of the chromatic polynomial of $C_n$
is infinitely log-concave if and only if $3\le n\le11$.
For $n=12$ its fifth iterate has a negative entry at degree two;
for $13\le n\le16$ its fourth iterate does; and for every $n\ge17$
its third iterate does.
\end{theorem}

In particular this family refutes Conjecture 21 of Amdeberhan and Moll
\cite{AM}, which proposes infinite log-concavity of the absolute coefficient
sequences of all chromatic polynomials. The same proposal appears as
Conjecture 13.1 of Amdeberhan \cite{A}. We use the explicit zero-extension
convention of Br\"and\'en \cite{B}; we do not silently replace it by
deleting endpoints at each iteration. No claim that $C_{12}$ is a smallest
counterexample among all finite graphs is made.

\section{The actual graph and its coefficients}
Deleting an edge of $C_n$ produces a path on $n$ vertices, and contracting
it produces $C_{n-1}$ when $n\ge4$. Deletion--contraction and the path
coloring count give
\[
 \chi_n(q)=q(q-1)^{n-1}-\chi_{n-1}(q),\qquad
 \chi_3(q)=q(q-1)(q-2).
\]
Induction proves
\begin{equation}\label{eq:poly}
 \chi_n(q)=(q-1)^n+(-1)^n(q-1).
\end{equation}
Consequently the full absolute coefficients, in ascending degree, are
\begin{equation}\label{eq:sequence}
 a_0=0,\qquad a_1=n-1,\qquad
 a_k=\binom nk\quad(2\le k\le n).
\end{equation}
The formal proof also verifies the genuine graph correspondence: it
constructs a bijection between proper colorings of Mathlib's
\texttt{SimpleGraph.cycleGraph n} with $q$ colors and closed walks of length
$n$ in the complete $q$-color graph. A trace computation gives
\eqref{eq:poly} for every $q\in\mathbb N$, including $q=0,1,2$.
Thus the formal sequence is not an assumed binomial surrogate for the
graph's chromatic polynomial.

\section{A uniform third-iterate obstruction}
Put $b=\mathcal La$ and $c=\mathcal Lb$. For $n\ge5$, substitution of
\eqref{eq:sequence} gives
\begin{align*}
b_1&=(n-1)^2,&
b_2&=\frac{n(n-1)^2(n+4)}{12},\\
b_3&=\frac{n^2(n-2)(n-1)^2(n+1)}{144},&
b_4&=\frac{n^2(n-3)(n-2)^2(n-1)^2(n+1)}{2880}.
\end{align*}
A second subtraction gives
\begin{align*}
c_1&=(n-1)^4,\\
c_2&=\frac{n^2(n-1)^4(n+2)}{16},\\
c_3&=\frac{n^3(n-2)^2(n-1)^4(n+1)(n^2+n+18)}{51840}.
\end{align*}
These equalities follow by replacing each binomial coefficient by its
product formula and collecting factors; all denominators are fixed positive
integers. Subtracting $c_1c_3$ from $c_2^2$ yields the exact identity
\begin{equation}\label{eq:third}
 (\mathcal L^3a)_2
 =-\frac{n^3(n-1)^8(n+4)}{103680}
       \bigl(2n^4-12n^3-327n^2-412n+36\bigr).
\end{equation}
With $n=u+17$, the final factor becomes
\[
 2u^4+124u^3+2529u^2+17370u+6615.
\]
It is strictly positive for every $u\ge0$. All other factors in
\eqref{eq:third}, apart from its minus sign, are positive for $n\ge17$.
This proves the entire infinite negative family in Theorem~\ref{thm:classification}
by one identity and a uniform sign argument.

For example, $C_{17}$ starts with
$(a_0,\ldots,a_5)=(0,16,136,680,2380,6188)$ and
$(b_0,\ldots,b_4)=(0,256,7616,138720,1456560)$.
Then $c_1=65536$, $c_2=22491136$, $c_3=8150077440$, and
\[
 (\mathcal L^3a)_2=22491136^2-65536\cdot8150077440
                 =-28272276537344.
\]

\section{The remaining negative cases}
Starting from \eqref{eq:sequence}, exact repeated integer subtraction gives
\begin{center}
\begin{tabular}{ccl}
\toprule
$n$ & iterate $r$ & $(\mathcal L^r a)_2$\\\midrule
12 & 5 & $-249621701601023742801969101519201265582387397744$\\
13 & 4 & $-3618341131654935620812800$\\
14 & 4 & $-199158562975246657489530096$\\
15 & 4 & $-2734560032125157358883149375$\\
16 & 4 & $-25982668618402950000000000000$\\\bottomrule
\end{tabular}
\end{center}
The particularly large entry for $n=12$ can be checked from its three
preceding neighbors:
\[
 (\mathcal L^4 a)_1=45949729863572161,\qquad
 (\mathcal L^4 a)_2=1505318939996586804452286,
\]
\[
 (\mathcal L^4 a)_3=54746933663864342589859620161140.
\]
The square of the middle number minus the product of the outer two is
the displayed fifth-iterate value. The companion certificate contains the
complete preceding sequences for every case, and the Lean declarations
\texttt{cycle12\_negative} through \texttt{cycle16\_negative} check the same
integer statements by kernel-checked finite arithmetic.

\section{An invariant proving infinitely many positive iterates}
\begin{lemma}\label{lem:invariant}
Suppose a nonnegative zero-extended finite sequence $v$ satisfies
$v_k^2\ge3v_{k-1}v_{k+1}$ at every index. Then $w=\mathcal Lv$ is
nonnegative and satisfies the same inequalities.
\end{lemma}
\begin{proof}
For every index,
\[
 \frac23v_k^2\le w_k\le v_k^2.
\]
Thus $w$ is nonnegative, and
\[
 w_{k-1}w_{k+1}\le v_{k-1}^2v_{k+1}^2\le\frac19v_k^4,
 \qquad w_k^2\ge\frac49v_k^4.
\]
It follows that $w_k^2\ge4w_{k-1}w_{k+1}\ge3w_{k-1}w_{k+1}$.
Zeros at the boundaries cause no division or exceptional case.
\end{proof}
If finitely many nonnegative initial iterates reach this invariant, induction
using Lemma~\ref{lem:invariant} proves nonnegativity of every subsequent
iterate. For each $3\le n\le11$ the following complete finite certificate
specifies such an iterate $j$. Its third column is the minimum of
$v_k^2-3v_{k-1}v_{k+1}$ over $2\le k\le n-1$, for $v=\mathcal L^j a$.
\begin{center}
\begin{tabular}{rrr}
\toprule
$n$ & $j$ & minimum exact slack\\\midrule
3&0&3\\4&1&28\\5&2&25825\\6&2&90846\\7&2&271656\\
8&3&711608924160\\9&3&4029124698225\\
10&3&19225238518750\\11&3&79738750726206\\\bottomrule
\end{tabular}
\end{center}
Each row is a directly specified finite integer computation: form
\eqref{eq:sequence}, iterate the displayed local rule $j$ times, check every
preceding entry is nonnegative, then check the listed minimum. Full
intermediate arrays are supplied in \texttt{exact\_certificates.json}.
At index one the slack is $v_1^2$ because $v_0=0$; at index zero and the
right boundary the assertions are automatic. All indices outside the
support remain zero under iteration. This verifies every index, and the
invariant proves the positive half of Theorem~\ref{thm:classification}.
Together with the negative certificates this completes its proof.

\section{Exact reproduction and limits}
Run \texttt{python3 verify\_cycle\_counterexample.py} from this paper's
directory. The standard-library checker recomputes the binomial arrays,
all finite certificates, deletion--contraction polynomial identities and
\eqref{eq:third} over exact rational polynomials. Finite calculation proves
only the finitely specified certificates; the infinite claims use the
invariant and the shifted positive polynomial above, both fully formalized.
The result classifies cycles, not all finite graphs or all chromatic
coefficient sequences. The earlier public $C_{17}$ counterexample and
written classification are retained and credited as predecessors of this
full graph-family formalization.

\section{Formal verification and reproduction}
The principal theorem has a complete Lean proof. The precise scope is:
Every genuine cycle graph $C_n$, n $\ge$ 3; full absolute coefficients, zero extension; infinitely log-concave iff n $\le$ 11. The graph/polynomial correspondence and all positive and negative cases are formalized.
The public proof source snapshot is fixed at commit
\nolinkurl{02af7db99db789fecb0d5f5f8f37b1a4e40762a0} in \url{https://github.com/mxym/math}.
The proof packages, with their own dependency locks, audit records and
reproduction instructions, are:
\begin{itemize}
\item \nolinkurl{formalizations/chromatic-cycles-all-n-lean}.
\end{itemize}
The main formal entry points are:
\begin{itemize}
\item \nolinkurl{ChromaticCycleAll.cycleGraph_chromatic_polynomial_classification}.
\item \nolinkurl{ChromaticCycleAll.every_actual_cycle_n_ge_17_fails}.
\end{itemize}
Lean is fixed to version 4.34.1, compiler commit
\nolinkurl{5045d0056413266e57c625dcd7c365b10e377c52}; Mathlib is fixed to
\nolinkurl{d13f23b723b8a846827a245b89c10fc7d3f11612}.
The audited mathematical roots depend only on the standard axioms
\texttt{propext}, \texttt{Classical.choice}, and \texttt{Quot.sound}.
The archived audits state their actual compilation and replay scopes;
publication of this typeset version does not imply a new execution of every
previous empty-kernel replay. The companion README maps each printed
principal result to its formal theorem and explains how to repeat verification.

\section*{Research and preparation disclosures}
This research, formalization and manuscript preparation used AI assistance.
The author is responsible for the manuscript. The principal claims are
supported by the readable proofs and the public Lean artifacts specified
above. Formal verification is not external human peer review; no such review
is claimed. No external funding was received. No worldwide priority or
journal acceptance is asserted.
\begin{thebibliography}{9}
\bibitem{AM} T. Amdeberhan and V. H. Moll, \emph{Infinite log-convexity}, Open Journal of Analytic Combinatorics \textbf{17} (2022), article 305, Conjecture 21. \url{https://combinatorialpress.com/article/ojac/vol17/305.pdf}.
\bibitem{A} T. Amdeberhan, \emph{Theorems, Problems and Conjectures}, arXiv:1207.4045v7 (2022), Conjecture 13.1. \url{https://arxiv.org/abs/1207.4045v7}.
\bibitem{B} P. Br\"and\'en, \emph{Iterated sequences and the geometry of zeros}, arXiv:0909.1927, especially the zero-extension convention. \url{https://arxiv.org/abs/0909.1927}.
\end{thebibliography}
\end{document}
